\section{Security Assessment and Risk Analysis}
\label{sec:security_assessment_risk_analysis}

Security assessment plays a fundamental role in secure software development by identifying vulnerabilities throughout the software lifecycle before they can be exploited in production environments. Rather than treating security as a separate activity performed after software implementation, modern secure software development practices advocate continuous security assessment integrated into development, deployment, and operational processes. The NIST Secure Software Development Framework (SSDF) recommends incorporating automated security verification throughout the software development lifecycle to reduce software vulnerabilities and improve software integrity \cite{NISTSSDF2022}. Similarly, the OWASP Software Assurance Maturity Model (SAMM) identifies continuous verification, security testing, and defect management as essential practices for establishing mature software security programs \cite{OWASPSAMM2024}.

A wide range of security assessment techniques has been proposed to identify vulnerabilities from different perspectives. Static Application Security Testing (SAST) analyzes source code without execution to detect programming flaws, insecure coding patterns, and potential vulnerabilities during software development. In contrast, Software Composition Analysis (SCA) focuses on identifying vulnerabilities introduced through third-party libraries and software dependencies by comparing dependency information against publicly disclosed vulnerability databases. Additional techniques, including infrastructure and configuration analysis, evaluate deployment environments and cloud infrastructure to identify security misconfigurations that may expose software systems to unnecessary risks. Because each technique targets different aspects of software security, they are generally considered complementary rather than interchangeable \cite{NISTSSDF2022, OWASPSAMM2024}.

As software systems become increasingly complex, organizations commonly employ multiple security assessment tools throughout the software lifecycle. Each tool produces findings using its own severity classification, confidence level, and reporting format, interpreting assessment results increasingly challenging. Consequently, modern vulnerability management extends beyond vulnerability detection to include vulnerability prioritization, enabling organizations to identify which security findings should be addressed first according to their potential impact and exploitation risk. NIST emphasizes that effective vulnerability management requires continuous identification, assessment, prioritization, remediation, and verification of security vulnerabilities throughout system operations rather than treating vulnerability scanning as an isolated activity \cite{NIST80040r4}.

Risk evaluation therefore has become an essential component of modern security assessment. The Common Vulnerability Scoring System (CVSS), maintained by the Forum of Incident Response and Security Teams (FIRST), provides a standardized approach for measuring the severity of software vulnerabilities based on exploitability, impact, and environmental factors \cite{FIRSTCVSSv40}. Although CVSS has become the de facto standard for vulnerability severity assessment, it primarily evaluates individual vulnerabilities rather than supporting organizational remediation decisions. To address this limitation, the Stakeholder-Specific Vulnerability Categorization (SSVC) framework incorporates additional decision factors, including exploitation status, mission impact, and operational context, to assist organizations in prioritizing vulnerability remediation according to their specific environments \cite{SSVC2023}.

Recent research has increasingly explored the application of machine learning techniques to vulnerability assessment and security risk prediction. Unlike traditional rule-based scoring systems, machine learning models are capable of learning relationships among multiple security features and predicting vulnerability severity or exploitation likelihood from historical data. Various supervised learning algorithms have been investigated for this purpose, including Logistic Regression, Decision Trees, Random Forests, Support Vector Machines, and Gradient Boosting models. These approaches have been applied to vulnerability prioritization, exploit prediction, and software defect prediction, demonstrating that machine learning can improve the consistency and scalability of security risk evaluation by combining heterogeneous security indicators into a unified predictive model \cite{Bozorgi2010BeyondHeuristics, Shin2013SoftwareMetrics, Rahman2019VulnerabilityPrediction}.

In practice, organizations often combine multiple security assessment techniques to obtain a more comprehensive understanding of software security. Static analysis, dependency analysis, infrastructure assessment, and compliance verification provide complementary perspectives on software risk, motivating the integration of heterogeneous security findings into unified security assessment workflows. Numerous tools have been developed to support these techniques, including Bandit for Python static analysis, Semgrep for pattern-based static analysis across multiple programming languages, and OWASP Dependency-Check for software dependency vulnerability analysis. While these tools improve vulnerability detection individually, effectively integrating their findings into a unified security evaluation process remains an important challenge for secure software operations. This challenge motivates the unified security assessment mechanism proposed in this research.